\par\medskip
\noindent\underline{Demostración de \sisi{2.5.11}{2.3.11}.}
\footnote{Para una demostración más elegante, véase 2.3.19 más abajo.}
Sea $\xi$ un elemento de $\Biext^1(P,Q;G)$, sea
\[
 \varphi_\xi:{}_nP\otimes{}_nQ\longrightarrow G
\]
el acoplamiento asociado, y sea
\[
 u_\xi:{}_nP\longrightarrow{}_n\uExt^1(nQ,G)
\]
la imagen de $\xi$ por (2.3.5). Hay que demostrar que $-u_\xi$ se
obtiene de $\varphi_\xi$ mediante la flecha (2.3.9), es decir, que es
el compuesto
\begin{equation}\tag{*}
 {}_nP
 \xrightarrow{\ \varphi_\xi\ }
 \uHom({}_nQ,G)
 \xrightarrow{\ (2.3.8)\ }
 {}_n\uExt^1(nQ,G).
\end{equation}
Para probarlo, podemos tomar una sección $x$ de $\,{}_nP$ y demostrar
que su imagen por $-u_\xi$ es igual a su imagen por el compuesto (*).
Si se desea, esto permitiría reducirnos al caso
$P=(\mathbb Z/n\mathbb Z)_T$, pero no necesitaremos esta reducción.

\medskip
\noindent a) Elijamos resoluciones planas $L_\bullet$ y $M_\bullet$,
\[
 0\longrightarrow L_1\longrightarrow L_0\longrightarrow0,
 \qquad
 0\longrightarrow M_1\longrightarrow M_0\longrightarrow0,
\]
de $P$ y $Q$, respectivamente.

%%%%%%%%%%  scan idx 255  |  folio 244  %%%%%%%%%%

Por ejemplo, se obtiene una resolución plana de $P$ escribiendo $P$
como cociente de un $\mathbb Z_T$-módulo plano $L_0$ y tomando como
$L_1$ el núcleo de $L_0\longrightarrow P$; este núcleo es plano
porque carece de torsión. Obtenemos el complejo
$L_\bullet\otimes M_\bullet$
\[
 0\longrightarrow L_1\otimes M_1
 \xrightarrow{d_2}
 L_1\otimes M_0+L_0\otimes M_1
 \xrightarrow{d_1}
 L_0\otimes M_0
 \longrightarrow0,
\]
cuyo haz $\underline H_1$ servirá como $\uTor_1(P,Q)$.

Sea $L'_0$ (respectivamente, $M'_0$) el subhaz de $L_0$
(respectivamente, de $M_0$) que es la imagen inversa de $\,{}_nP$
(respectivamente, de $\,{}_nQ$). Así,
\[
 0\longrightarrow L_1\longrightarrow L'_0\longrightarrow0
 \qquad
 \left(\text{respectivamente, }
 0\longrightarrow M_1\longrightarrow M'_0\longrightarrow0\right)
\]
es una resolución plana de $\,{}_nP$ (respectivamente, de
$\,{}_nQ$). El homomorfismo
\[
 {}_nP\otimes{}_nQ
 \longrightarrow
 \uTor_1(P,Q)=\underline H_1(L_\bullet\otimes M_\bullet)
\]
puede describirse entonces, pasando al cociente, a partir del
homomorfismo
\[
 L'_0\otimes M'_0
 \longrightarrow
 \underline H_1(L_\bullet\otimes M_\bullet)
\]
dado por
\[
 x_0\otimes y_0
 \longmapsto
 \text{la clase del $1$-ciclo }
 [nx_0]\otimes y_0-x_0\otimes[ny_0].
\]
Los corchetes indican que el elemento considerado se ve como una
sección del subhaz $L_1$ de $L_0$ (respectivamente, de $M_1$ en
$M_0$). Esta descripción se deduce del cálculo efectuado en la
demostración de 2.1.10.

\medskip
\noindent b) Tomemos una resolución inyectiva
\[
 C^\bullet=\bigl(0\longrightarrow C^0
 \xrightarrow{d^0}C^1\longrightarrow\cdots\bigr)
\]
de $G$, e identifiquemos
\[
 \Biext^1(P,Q;G)
 \simeq
 \Ext^1\!\left(P\mathbin{\overset{\mathbb L}{\otimes}}Q,G\right)
 =H^1\!\left(
   \Hom^{\boldsymbol\cdot}
   (L_\bullet\otimes M_\bullet,C^{\boldsymbol\cdot})
 \right).
\]
El elemento $\xi$ puede representarse, por tanto, mediante un par de
homomorfismos $f_1,f_0$ que dan el diagrama anticonmutativo

%%%%%%%%%%  scan idx 256  |  folio 245  %%%%%%%%%%

\begin{equation}\tag{2.3.12}
\begin{tikzcd}[column sep=huge,row sep=huge]
 L_1\otimes M_0+L_0\otimes M_1
   \arrow[r,"d_1"]
   \arrow[d,"f_1"']
   \arrow[rd,phantom,"-" description]&
 L_0\otimes M_0
   \arrow[d,"f_0"']\\
 C^0\arrow[r,"d^0"]&C^1.
\end{tikzcd}
\end{equation}
Por definición, el acoplamiento
\[
 \varphi_\xi:{}_nP\otimes{}_nQ\longrightarrow G
\]
se obtiene, pasando al cociente, a partir del homomorfismo
\[
 \varphi:L'_0\otimes M'_0\longrightarrow G
\]
dado por
\begin{equation}\tag{2.3.13}
 \varphi(x_0\otimes y_0)
 =f_1\bigl([nx_0]\otimes y_0-x_0\otimes[ny_0]\bigr)
 \in G=\ker(C^0\longrightarrow C^1).
\end{equation}
Como la afirmación 2.3.11 es local sobre $T$, podemos suponer que la
sección dada $x$ de $\,{}_nP$ se levanta a una sección $x_0$ de
$L'_0$. Considerando $y_0$ variable, la fórmula precedente define
entonces un homomorfismo
\begin{equation}\tag{2.3.14}
 \varphi':M'_0\longrightarrow G,
 \qquad
 y_0\longmapsto
 f_1\bigl([nx_0]\otimes y_0
 \sisi{\longrightarrow}{-}x_0\otimes[ny_0]\bigr),
\end{equation}
que pasa al cociente para dar
\[
 {}_nQ\longrightarrow G,
 \qquad
 y\longmapsto\varphi_\xi(x,y),
\]
el homomorfismo asociado a $\varphi_\xi$ y $x$.
